\label{sec:task_formulation}
{This work introduces a new task formulation, \emph{retrieval with instructions} (Figure~\ref{fig:teaser} bottom right).
}
{We are given} a large collection of $N$ documents $\mathcal{D} = \{\boldsymbol{d}_1, \ldots, \boldsymbol{d}_N\}$, a search task instruction $\boldsymbol{t}$ and a query $\boldsymbol{q}$. 
The problem of retrieval with instructions is to find a document $\boldsymbol{d} \in \mathcal{D}$ that is relevant to $\boldsymbol{q}$ according to the instruction $\boldsymbol{t}$.
Compared to the standard retrieval problem setting (e.g., Figure~\ref{fig:teaser} bottom left), the difference is the explicit definition of \emph{relevance} in the instruction~$\boldsymbol{t}$ as additional input to the system, which can be diverse and may not be fully defined by the query only~\cite{ruthven_lalmas_2003}. 
Thus, retrieval with instructions lets us build retrieval systems that are very general and task-aware---changing their relevance measure by attending to the instruction.

{
This new formulation brings both new research challenges and opportunities. For instance, a retriever is now required to modify its search behavior according to the instructions. 
On the plus side, different datasets can be naturally grouped to train a single retriever, yielding benefits from cross-task interdependence, which has been observed in instruction tuning of LLMs.
% , and can adapt to unseen tasks via instructions. 
Compared to training a retriever on multiple tasks with task identifiers~\cite{maillard-etal-2021-multi}, instructions provide greater flexibility and enable zero-shot transfer via natural language instructions.
% when the evaluation task does not resemble any of the training tasks. 
Furthermore, a multi-task instruction-following retriever obviates the need to host multiple task-specific retrievers.
}

{
As noted previously, multi-task training with instructions has not been studied in the area of retrieval due to the lack of resources and dedicated models. To facilitate the research on retrieval with instructions, we  introduce the first large-scale retrieval benchmark with expert-written annotations (Section~\ref{sec:dataaset}) and subsequently the multi-task instruction-following retrievers (Section~\ref{sec:method})}. 


